% ==============================================================================
% sec:datova_analyza --- metodika zpracování dat
% ==============================================================================

Datová část práce je vystavěna jako reprodukovatelný řetězec -- zdrojová data se
stahují a~zpracovávají programaticky ve skriptech v~jazyce Python,
a~výstupy se zapisují do~grafů a~tabulek vkládaných do~\LaTeX{}u. Doprovodný text
slouží k~vysvětlení, jak byl ukazatel sestaven, jaký má rozsah a~jaký význam má
pro argument práce. Pokud není uvedeno jinak, mapové a~časové grafy zachycují
nejnovější společný rok dostupný ve zdrojových datech a~chybějící hodnoty nejsou
nahrazovány výpočtem nebo odhadem. Bližší dokumentace programatického zpracování je dostupná v~příloze~1.\par

\paragraph{Harmonizovaná evropská srovnání}\label{par:data_zdroje} vycházejí primárně z~Eurostatu a~z~databází
\acs{OECD}, zejména \acs{ICTWSS} pro odborovou organizovanost a~pokrytí
\aclins{KS} a~\acs{LMP} pro výdaje na~\aclacc{APZ}. Národní mzdové analýzy pro
\ac{ČR} využívají \acs{ISPV}, zatímco vyjednané mzdové a~pracovněprávní
parametry kolektivních smluv vycházejí z~\acs{IPP}.\par

\paragraph{Chybějící hodnoty}\label{par:missing} se zobrazují jako
\texttt{--}, v~mapách jako šedá území. Důvodem může být pozdní aktualizace,
nízká dostupnost dat nebo skutečnost, že zdroj nepokrývá všechny země
stejně. Zvláštní pozornost vyžadují databáze \acgen{OECD}, které jsou vymezeny
členstvím v~\acloc{OECD}, nikoli členstvím v~\acloc{EU}. Proto se v~datech
\acs{ICTWSS} a~\acs{LMP} objevují i~země mimo \acacc{EU}, zatímco některé
členské státy mají neúplné datové řady. Datové sady Eurostatu
většinou pokrývají všech 27 členských států \acgen{EU}, pokud státy data vykazují.\par
Některé použité datové sady jsou sestavovány méně než posledních 10~let, dostupnost
dat pro analýzy trhu práce se ale výrazně zlepšila, a~lze v~nich zachytit důležité trendy.
U~zemí s~datovou mezerou mohla být v~průřezových analýzách použita starší data, i~tato odchylka se uvádí 
u~těch výstupů, které jsou dotčeny.\par

\paragraph{Přepočet kupní síly}\label{par:pps_metodika} je pro mezistátní srovnání příjmů, nákladů práce a~produktivity nutný, protože 
nominální hodnota nezohledňuje rozdíly v~cenové hladině mezi zeměmi. Eurostat k~tomu publikuje 
\acf{PPP} a~odvozený agregát \acf{PPS}. Převod nominální hodnoty vyjádřené
v~národní měně na~hodnotu podle kupní síly (\acs{PPS}) probíhá dělením
příslušnou paritou kupní síly \(\acs{PPP}_{\acs{idx-i}}\):
\begin{equation}
  \acs{var-x}^{\acs{PPS}}_{\acs{idx-i}} = \frac{\acs{var-x}_{\acs{idx-i}}}{\acs{PPP}_{\acs{idx-i}}}.
  \label{eq:pps_konverze}
\end{equation}
Pakliže jsou hodnoty v~databázi dostupné pouze v~eurech, používá se konzistentní převod přes
index srovnatelné cenové hladiny \(\acs{PLI}_{\acs{idx-i}}\) (Price Level Index) z~datasetu \texttt{prc\_ppp\_ind}:
\begin{equation}
  \acs{var-x}^{\acs{PPS}}_{\acs{idx-i}} = \frac{\acs{var-x}^{\acs{EUR}}_{\acs{idx-i}}}{\acs{PLI}_{\acs{idx-i}}/100},
  \label{eq:lci_eur_pps}
\end{equation}
kde \(\acs{var-x}^{\acs{EUR}}_{\acs{idx-i}}\) představuje nominální hodnotu v~eurech. Pro analýzy nákladů práce
z~datasetu \texttt{lc\_lci\_lev} se uplatňuje tento převod za použití cenové hladiny \(\acs{PLI}\) odpovídající \acloc{HDP} (\texttt{PLI\_EU27\_2020.GDP}).
Pokud poslední rok nákladových dat ještě nemá publikovaný index cenové hladiny, využívá se nejbližší dostupný předchozí rok. Tento postup
je jednotně použit napříč celou prací tam, kde se porovnávají paritně očištěné mzdové nebo nákladové ukazatele mezi zeměmi.\par

\paragraph{Odvozené hodinové ukazatele}\label{par:hodinove_ukazatele} vstupují do~korelační a~komparativní části jako hodinový čistý příjem v~\aclloc{PPS},
protože roční čistý příjem samotný nepostihuje rozdíly v~délce skutečné odpracované
doby. Roční čistý příjem standardizované osoby \acs{idx-AW100} se dělí
odhadovaným ročním pracovním fondem, který vzniká násobením skutečné týdenní
odpracované doby průměrným počtem týdnů v~roce:
\begin{equation}
  \acs{var-w}^{\acs{PPS}}_{\acs{idx-n},\acs{idx-h},\acs{idx-AW100}}
    = \frac{\acs{var-y}^{\acs{PPS}}_{\acs{idx-n},\acs{idx-AW100},\acs{idx-a}}}
           {\acs{var-H}_{\acs{idx-w}} \cdot 52{,}18}.
  \label{eq:pps_hourly}
\end{equation}
\par

\paragraph{Indexové řady a~srovnávací osa}\label{par:normalizace} jsou používány tam, kde je podstatnější relativní pozice a~trend konvergence než
absolutní nominální úroveň. U~\aclgen{HDP} na obyvatele, produktivity práce,
cenové hladiny a~některých příjmových ukazatelů je referenční hodnotou průměr
\aclgen{EU} nastavený na~100. Stabilní srovnávací šestici tvoří \ac{ČR}, Rakousko, Německo, Dánsko,
Polsko a~Slovensko; mapy a~korelační analýzy pracují podle dostupnosti dat se
širším evropským panelem. U~kompozitního modelu (ternární analýzy)
jsou dílčí proměnné převáděny metodou min-max na společnou škálu \SIrange{0}{100}{}:
\begin{equation}
  \acs{var-z}^{\ast}_{\acs{idx-i},\acs{idx-k}}
    = \frac{\acs{var-x}_{\acs{idx-i},\acs{idx-k}} - \min\limits_{\acs{idx-k}\in\acs{geo-EU27}} \acs{var-x}_{\acs{idx-i},\acs{idx-k}}}
           {\max\limits_{\acs{idx-k}\in\acs{geo-EU27}} \acs{var-x}_{\acs{idx-i},\acs{idx-k}} - \min\limits_{\acs{idx-k}\in\acs{geo-EU27}} \acs{var-x}_{\acs{idx-i},\acs{idx-k}}}
      \cdot 100,
  \label{eq:minmax_general}
\end{equation}
kde \(\acs{var-x}_{\acs{idx-i},\acs{idx-k}}\) je hodnota indikátoru \(\acs{idx-i}\) pro zemi \(\acs{idx-k}\) a~jmenovatel
vyjadřuje rozpětí napříč panelovou skupinou \(\acs{geo-EU27}\). U~indikátorů,
u~nichž vyšší hodnota znamená zhoršení výsledku, se používá inverzní transformace:
\begin{equation}
  \tilde{\acs{var-z}}^{\ast}_{\acs{idx-i},\acs{idx-k}} = 100 - \acs{var-z}^{\ast}_{\acs{idx-i},\acs{idx-k}}.
  \label{eq:minmax_inverse_general}
\end{equation}
Po normalizaci jsou jednotlivé složky rozhraní mezi sebou číselně porovnatelné a~lze je agregovat
váženým průměrem do výsledných kompozitních os.\par

\paragraph{Daňový klín standardizovaného zaměstnance}\label{par:tax_distribution} je vyjádřen jako podíl daně a~povinných
odvodů na celkových nákladech práce:
\begin{equation}
  \acs{var-tau}
    = \frac{\acs{var-c}_{\acs{SP},\acs{idx-z}}
            + \acs{var-c}_{\acs{SP},\acs{idx-Z}}
            + \acs{var-c}_{\acs{DPFO}}}
           {\acs{var-w}_{\acs{idx-b}} + \acs{var-c}_{\acs{SP},\acs{idx-Z}}}.
  \label{eq:tax_wedge}
\end{equation}
U~modelů \acs{OSVČ} je stejná logika odvodového zatížení aplikována na příjmový nebo ziskový základ
podle zvoleného režimu. Při výdajovém paušálu se základ daně odvozuje z~příjmu
po odečtení paušálních výdajů, zatímco u~paušální daně je platba v~daném pásmu
pevná.\par

\paragraph{Modelové výpočty pro \acacc{ČR}}\label{par:cz_model} vycházejí u~čistého příjmu, sociálního a~zdravotního pojištění, daňového klínu
a~starobního důchodu z~legislativního rámce \acgen{ČR} platného pro
rok~2026: ze zákona o~daních z~příjmů, zákoníku práce, zákona o~důchodovém
pojištění, důchodové reformy, zákonů o~sociálním a~zdravotním pojištění
a~prováděcího nařízení vlády.
Výsledky odpovídají základnímu pracovníkovi bez slev na \acloc{DPFO} mimo základní slevu na
poplatníka. U~zaměstnance je osobní vyměřovací základ odvozen z~hrubé mzdy,
zatímco vodorovná osa modelů pracuje s~celkovým nákladem zaměstnavatele. U~\acs{OSVČ}
jsou rozlišeny standardní režim a~paušální daň, protože obě konstrukce odlišně
převádějí stejné prostředky do~odvodů, čistého příjmu a~budoucího starobního
důchodu.\par

Důchodové modely předpokládají pojistnou dobu 40~let a~počítají základní
i~procentní výměru podle redukčních hranic platných pro rok~2026.\par
Čisté příjmy
jsou hodnotami před exekučními a~insolvenčními srážkami. Tento limit je pro
interpretaci nízkopříjmových skupin podstatný, neboť stejná srážka
představuje u~nižšího příjmu vyšší podíl disponibilního příjmu.~\cite{PAQ_ChudobaCEEU_2025}\par

\paragraph{Příjmová nerovnost}\label{par:gini_y_sec} je v~evropských srovnáních zachycena Giniho koeficientem
ekvivalizovaného disponibilního příjmu domácností:
\begin{equation}
  \acs{var-g}_{\acs{idx-y}}
    = \frac{2\sum_{\acs{idx-i}=1}^{\acs{var-N}}
              \acs{idx-i}\;\acs{var-y}_{\acs{idx-i}}}
           {\acs{var-N}\sum_{\acs{idx-i}=1}^{\acs{var-N}}
              \acs{var-y}_{\acs{idx-i}}}
    - \frac{\acs{var-N}+1}{\acs{var-N}},
  \label{eq:gini_y}
\end{equation}
kde se předpokládá uspořádání příjmů jednotlivců neklesajícím způsobem \(\acs{var-y}_1 \le \acs{var-y}_2 \le \dots \le \acs{var-y}_{\acs{var-N}}\). Grafy pracující s~decilovými hranicemi příjmů nebo mezd používají log-normální
interpolaci nad publikovanými kvantilovými body. Jde o~vyhlazení rozdělení pro vizualizaci,
nikoli o~dodatečný statistický odhad individuálních dat.\par

\paragraph{Metodika ternárního modelu}\label{par:ternary_metodika} je v~práci jednotná napříč všemi výstupy, které pracují s~rovnováhou
priorit zaměstnanců, zaměstnavatelů a~vlády. Pro každého aktéra je nejprve
sestaven vážený kompozitní index z~normalizovaných dílčích indikátorů:
\begin{equation}
  \acs{var-S}_{\acs{idx-actor},\acs{idx-k}} = \sum_{\acs{idx-i}=1}^{\acs{var-m_ind}_{\acs{idx-actor}}} \acs{var-omega}_{\acs{idx-actor},\acs{idx-i}}\, \acs{var-z}^{\ast}_{\acs{idx-actor},\acs{idx-i},\acs{idx-k}},
  \qquad
  \sum_{\acs{idx-i}=1}^{\acs{var-m_ind}_{\acs{idx-actor}}} \acs{var-omega}_{\acs{idx-actor},\acs{idx-i}} = 1,
  \label{eq:ternary_axis_score}
\end{equation}
kde \(\acs{var-S}_{\acs{idx-actor},\acs{idx-k}}\) je skóre aktéra \(\acs{idx-actor}\) v~zemi \(\acs{idx-k}\), \(\acs{var-z}^{\ast}_{\acs{idx-actor},\acs{idx-i},\acs{idx-k}}\) je min-max
normalizovaná hodnota dílčí proměnné a~\(\acs{var-omega}_{\acs{idx-actor},\acs{idx-i}}\) její váha. Následně se
provádí normalizace na~součet \SI{100}{\percent}:
\begin{equation}
  \acs{var-P_share}_{\acs{idx-actor},\acs{idx-k}} = \frac{\acs{var-S}_{\acs{idx-actor},\acs{idx-k}}}{\sum_{\acs{idx-actor}\in\{\mathrm{zam.},\mathrm{firmy},\mathrm{vlada}\}} \acs{var-S}_{\acs{idx-actor},\acs{idx-k}}}\cdot 100,
  \label{eq:ternary_closure}
\end{equation}
takže výsledné podíly \(\acs{var-P_share}_{\acs{idx-actor},\acs{idx-k}}\) jsou mezi zeměmi přímo srovnatelné.\par

Pro nesektorový ukazatel míry volných pracovních míst je použit sektorový filtr (bez zemědělství a~veřejné správy), při nedostupnosti nahrazován v~pořadí 
\texttt{B-S\_X\_O} \(\rightarrow\) \texttt{TOTAL} \(\rightarrow\)
\texttt{B-S} \(\rightarrow\) \texttt{B-O}. Primárním zdrojem je
\texttt{jvs\_a\_r21} s~preferencí tříletého průměru (\texttt{AVG\_3Y}).
Odchylka je uvedena v~textu přímo u~příslušného obrázku nebo tabulky.\par

Proměnná sociálního smíru je konstruována jako semikvantitativní expertní
skóre dlouhodobé provozní spolehlivosti trhu práce z~pohledu zaměstnavatelů.
Ztracených dnů práce kvůli průmyslovým sporům (na~\SI{1000}{\person}) ze statistik \acs{MOP}
slouží jako kvantitativní vstup, ale výsledné přiřazení probíhá s~ohledem na nízkou spolehlivost dat v~pěti fixních úrovních.\par

\paragraph{Korelační aparát}\label{par:korelace_math} používá jako vysvětlující proměnnou pokrytí \acpins{KS}. Výstupními
proměnnými jsou skutečná týdenní odpracovaná doba, genderový rozdíl v~odměňování v~\%,
čistý hodinový příjem a~hodinová produktivita práce v~\aclloc{PPS}. Doplňkový
pár tvoří odborová organizovanost a~Giniho koeficient disponibilního příjmu.
Síla vztahu je měřena Pearsonovým korelačním koeficientem
\begin{equation}
  \acs{var-r}_{\acs{idx-x},\acs{idx-y}}
  = \frac{\displaystyle\sum_{\acs{idx-i}=1}^{\acs{var-n}}
            (\acs{var-x}_{\acs{idx-i}} - \bar{\acs{var-x}})
            (\acs{var-y}_{\acs{idx-i}} - \bar{\acs{var-y}})}
         {\displaystyle\sqrt{
            \sum_{\acs{idx-i}=1}^{\acs{var-n}}(\acs{var-x}_{\acs{idx-i}}-\bar{\acs{var-x}})^2
            \sum_{\acs{idx-i}=1}^{\acs{var-n}}(\acs{var-y}_{\acs{idx-i}}-\bar{\acs{var-y}})^2}},
  \label{eq:pearson_r}
\end{equation}
a~Spearmanovým pořadovým koeficientem
\begin{equation}
  \acs{var-rho} = \acs{var-r}_{\acs{idx-x},\acs{idx-y}}\!\left(\operatorname{R}(\acs{var-x}),\operatorname{R}(\acs{var-y})\right).
  \label{eq:spearman_rho}
\end{equation}
Pro modelový odhad změny při hypotetickém růstu pokrytí \aclins{KS} je použita
prostá lineární regrese
\begin{equation}
  \hat{\acs{var-y}}_{\acs{idx-i}} = \hat{\acs{var-alpha}} + \hat{\acs{var-beta}}\,\acs{var-x}_{\acs{idx-i}},
  \qquad
  \Delta\hat{\acs{var-y}} = \hat{\acs{var-beta}}\,\Delta \acs{var-x}.
  \label{eq:ols}
\end{equation}
Koeficienty se interpretují popisně. Panel zahrnuje opakovaná pozorování týchž
zemí v~čase, a~nejde tedy o~náhodný výběr nezávislých jednotek; korelace
není používána jako důkaz kauzality, nýbrž jako indikace síly asociace,
kterou je třeba zasadit do~výkladu pravidel a~mechanismů trhu práce.\par
